%%%%%%%%%%%%%%%%%%%%%%%%%%%%%%%%%%%%%%%%%%%%%%%%%%%%%%%%%%%%%%%%%%%%%%%%%%%%%%%%
%                                                                              %
% BlackArch Linux Guide                                                        %
%                                                                              %
%%%%%%%%%%%%%%%%%%%%%%%%%%%%%%%%%%%%%%%%%%%%%%%%%%%%%%%%%%%%%%%%%%%%%%%%%%%%%%%%


%%% BEGIN %%%
\documentclass[a4paper, oneside, 11pt]{book}


%%% INCLUDES %%%
\renewcommand{\familydefault}{\sfdefault}
\usepackage{verbatim}
\usepackage[utf8]{inputenc}
\usepackage{geometry}
\usepackage{array}
\usepackage{graphicx}
\usepackage{supertabular}
\usepackage{verbatim}
\usepackage{pstricks}
\usepackage{fancyhdr}
\usepackage{fancyvrb}
\usepackage{tocloft}
\usepackage{color}
\usepackage{enumerate}
\usepackage[french]{babel}
\usepackage[
pdftitle={BlackArch Linux, The BlackArch Linux Guide},
pdfsubject={BlackArch Linux, The BlackArch Linux Guide},
pdfauthor={BlackArch Linux, BlackArch Linux},
pdfkeywords={BlackArch Linux, Penetration Testing, Security, Hacking, Linux},
pdfcenterwindow=true
]{hyperref}


%%% LAYOUT %%%
\setlength{\parindent}{0em}
\setlength{\parskip}{1.5ex plus0.5ex minus0.5ex}
\geometry{left=2.5cm, textwidth=16cm, top=3cm, textheight=25cm, bottom=3cm}
\widowpenalty=2000
\clubpenalty=1000
\frenchspacing
\sloppy
\pagecolor{black}
\color{white}
\hypersetup{pdftex=true, colorlinks=true, breaklinks=true, linkcolor=red,
menucolor=red, pagecolor=red, urlcolor=red}
\setcounter{tocdepth}{10}
\setcounter{secnumdepth}{10}


%%% HEADER / FOOTER %%%
\setlength{\headheight}{1.0cm}
\setlength{\headsep}{1.0cm}
\lhead{{\includegraphics[width=1cm,height=1cm]{logo.png}}}
\rhead{The BlackArch Linux Guide}
%\lfoot{\small{foo}}
%\rfoot{\small{bar}}


%%% OWN MACROS %%%
% put own macros here, like renewcommand newcommand etc.


%%% DOCUMENT %%%
\begin{document}
\pagestyle{empty}
\begin{center}
\begin{figure}[htbp]
\centering
\vspace{1cm}
\Huge{\textbf{BlackArch \color{red}Linux}}\\
\vspace{2cm}
\includegraphics[width=8cm]{logo.png}
\label{fig:logo}
\end{figure}
\vspace{1cm}
\Huge{\textbf{Guide de BlackArch}}\\
\vspace{1cm}
\Large{\color{red}http://www.blackarch.org/}\\
\vspace{0.5cm}
\Large{\today}
\end{center}
\newpage
\tableofcontents
\newpage
\pagestyle{fancy}

%%%%%%%%%%%%%%%%%%%%%%%%%%%%%%%%%%%%%%%%%%%%%%%%%%%%%%%%%%%%%%%%%%%%%%%%%%%%%%%%

\chapter{Introduction}

\section{Qu'est-ce que BlackArch Linux?}
\href{http://www.blackarch.org}{BlackArch Linux} est une légère expansion d'Arch 
Linux destinée aux professionnels de la sécurité informatique.
\\\\
L'ensemble d'outils est distribué tel un  
\href{https://wiki.archlinux.org/index.php/Unofficial\_User\_Repositories}
{dépôt non-officiel} d'Arch Linux, vous permettant d'installer les composantes de
BlackArch sur une installation d'Arch Linux existante. Les paquets peuvent y être
installés individuellement ou par catégorie.
\\\\
Le dépot logiciel de Black Arch contient plus de 650 outils et ce nombre augmente
sans cesse. Tous ces outils sont minutieusement testés avant d'être ajouté aux
dépots, afin de garantir leurs bon fonctionnement.

%\section{The story of BlackArch Linux}
%foo bar

%\section{Supported platforms}
%foo bar

\section{S'impliquer}
Il est possible d'entrer en contact avec l'équipe Black Arch à l'aide de ces
informations:
\\\\
Web: \url{http://www.blackarch.org/}
\\\\
Courriel: \href{mailto:blackarchlinux@gmail.com}{blackarchlinux@gmail.com}
\\\\
IRC: \url{irc://irc.freenode.net/blackarch}

%%%%%%%%%%%%%%%%%%%%%%%%%%%%%%%%%%%%%%%%%%%%%%%%%%%%%%%%%%%%%%%%%%%%%%%%%%%%%%%%

\chapter{Guide d'utilisation}

\section{Installation}
La section suivante démontre comment configurer le dépôt de paquets de Black 
Arch ainsi que la procédure à suivre afin d'installer des logiciels à partir de 
celle-ci. BlackArch supporte l'installation de fichiers binaires précompilés 
ainsi que l'installation à l'aide du code source original des paquets.
\\\\
BlackArch est compatible avec une installation de base d'Arch Linux. Il s'agira 
d'un dépôt non-officiel de paquets. Si vous désirez une image ISO à la place,
visitez la section \href{http://www.blackarch.org/download.html#iso}{Live ISO}.
\\\\

\subsection{Installer le dépôt}
Il suffit de 6 étapes afin d'installer le dépot de BlackArch sur une installation
existante de Arch Linux. Vous devez suivre la recette dans l'ordre.
You must follow the instuctions in order. N'ajoutez pas \textbf{blackarch} 
à votre fichier \textit{pacman.conf} sans avoir complété les étapes 0 à 2.
\begin{enumerate} \itemsep4pt
\item Si vous avez déjà installé BlackArch et qu'il y existe déjà une entrée
\textbf{[blackarch]} dans \textit{/etc/pacman.conf}, vous devrez la supprimer ou
la mettre en commentaire, et ensuite exécuter la commande \textit{pacman -Syy}.
\item Exécutez les commandes suivantes en tant que root. Ces commande
téléchargeront et installeront les clés nécessaire pour garantir l'authenticité
des paquets via une signature digitale.
{\small
\color{gray}
\begin{verbatim}
wget -q http://blackarch.org/keyring/blackarch-keyring.pkg.tar.xz{,.sig}
gpg --keyserver hkp://pgp.mit.edu --recv 4345771566D76038C7FEB43863EC0ADBEA87E4E3
gpg --keyserver-o no-auto-key-retrieve --with-f blackarch-keyring.pkg.tar.xz.sig
pacman-key --init
rm blackarch-keyring.pkg.tar.xz.sig
pacman --noc -U blackarch-keyring.pkg.tar.xz
\end{verbatim}
}
\item Si désiré, vérifiez l'authenticité des clés autant de sources que vous 
désirez afin de vous convaincre de leurs validités.
\item Ajoutez les lignes suivantes à votre fichier \textit{/etc/pacman.conf}:
{\small
\color{gray}
\begin{verbatim}
[blackarch]
Server = <site_miroir>/$repo/os/$arch
\end{verbatim}
}
Replacez \textit{\textless site\_mirroir\textgreater} par le miroir officiel
de votre choix. La liste des miroirs est disponible sur
\href{http://www.blackarch.org/}{notre site web}.

\item Exécutez:
{\small
\color{gray}
\begin{verbatim}
# pacman -Syyu
\end{verbatim}
}
\end{enumerate}

\subsection{Installation de paquets}
Vous pouvez désormais installer des outils directement du dépôt de BlackArch
\begin{enumerate}
\item Afin d'obtenir la liste de tous les outils disponibles exécutez
{\small
\color{gray}
\begin{verbatim}
# pacman -Sgg | grep blackarch | cut -d' ' -f2 | sort -u
\end{verbatim}
}
\item Afin d'installer tous les outils, exécutez
{\small
\color{gray}
\begin{verbatim}
# pacman -S blackarch
\end{verbatim}
}
\item Afin d'installer une catégorie complète d'outils, exécutez
{\small
\color{gray}
\begin{verbatim}
# pacman -S blackarch-<nom de la catégorie>
\end{verbatim}
}
\item Afin d'obtenir la liste des catégories de paquets disponibles, exécutez
{\small
\color{gray}
\begin{verbatim}
# pacman -Sg | grep blackarch
\end{verbatim}
}
\end{enumerate}

\subsection{Installation des paquets depuis le code source}
Il existe une méthode alternative à l'installation de paquets sous forme binaire.
BlackArch vous offre la possibilité d'obtenir le code source directement du 
dépôt officiel du paquet désiré, et BlackArch le compilera pour vous. Vous 
trouverez les fichiers PKGBUILD sur notre 
\href{https://github.com/BlackArch/blackarch/tree/master/packages}{github}. Afin
d'obtenir le dépôt en entier, vous pourrez utiliser l'outil 
\href{https://github.com/BlackArch/blackman}{blackman}.
\begin{itemize} \itemsep10pt
\item Vous devez premièrement installer blackman. Si vous avez déjà configuré
le dépôt de paquets de BlackArch sur votre système, tel qu'indiqué au début de 
ce document, vous vous pouvez installer blackman à l'aide de pacman à l'aide de
cette commande:
{\small
\color{gray}
\begin{verbatim}
pacman -S blackman
\end{verbatim}
}
\item Vous pouvez aussi installer blackman directement du code source:
{\small
\color{gray}
\begin{verbatim}
mkdir blackman
cd blackman
wget https://raw2.github.com/BlackArch/blackarch/master/packages/blackman/PKGBUILD
# Make sure the PKGBUILD has not been maliciously tampered with.
makepkg -s
\end{verbatim}
}
\item Afin de télécharger et compiler un paquet:
{\small
\color{gray}
\begin{verbatim}
$ sudo blackman -i package
\end{verbatim}
}
\item Afin de télécharger, compiler, et installer une catégorie entière:
{\small
\color{gray}
\begin{verbatim}
$ sudo blackman -g group
\end{verbatim}
}
\item Afin de télécharger, compiler, et installer tous les outils présents
dans le dépôt de paquets de BlackArch:
{\small
\color{gray}
\begin{verbatim}
$ sudo blackman -a
\end{verbatim}
}
\item Afin d'obtenir la liste de toutes les catégories de paquets de BlackArch:
{\small
\color{gray}
\begin{verbatim}
$ blackman -l
\end{verbatim}
}
\item Si vous voulez ensuite obtenir la liste de tous les paquets compris dans
une catégorie particulière, vous pouvez obtenir la liste de ceux-ci ainsi qu'une
courte description pour chaque paquets à l'aide de la commande suivante:
{\small
\color{gray}
\begin{verbatim}
$ blackman -p <nom de la catégorie>
\end{verbatim}
}
\end{itemize}

\subsection{Installation d'un LiveCD, netinstall, ou d'Arch Linux}
Vous pouvez installer Black Arch d'un de nos images ISO, que ce soit un LiveCD
ou un netcd. Vous les trouverez sur 
\href{http://www.blackarch.org/download.html#iso}{notre site}. Les étapes 
suivantes sont nécéssaire afin de créer une image bootable:
\begin{itemize}
\item Installez le paquet blackarch-install-scripts:
{\small
\color{gray}
\begin{verbatim}
$ sudo pacman -S blackarch-install-scripts
\end{verbatim}
}
\item Ensuite exécutez:
{\small
\color{gray}
\begin{verbatim}
$ sudo blackarch-install
\end{verbatim}
}
\end{itemize}

%%%%%%%%%%%%%%%%%%%%%%%%%%%%%%%%%%%%%%%%%%%%%%%%%%%%%%%%%%%%%%%%%%%%%%%%%%%%%%%%


\chapter{Guide de développement}

\section{Le Arch Build System et les dépôts de paquets}

Les fichiers PKGBUILD sont de simples scripts permettant d'installer et de 
compiler et installer un programme sur plusieurs environnements différents. 
Chaque PKGBUILD indique à la commande makepkg(1) comment créer un paquet valide
qui pourra ensuite être installé. Les fichiers PKGBUILD sont écrits en bash.

Pour plus d'information, il est suggéré de lire ces pages:
\begin{itemize}
\item \href{http://wiki.archlinux.fr/Standard_paquetage}{Arch Wiki FR:
Standard paquetage}
\item \href{http://wiki.archlinux.fr/Makepkg}{Arch Wiki FR: Makepkg}
\item \href{http://wiki.archlinux.fr/PKGBUILD}{Arch Wiki FR: PKGBUILD}
\end{itemize}

\section{Le standard de paquetage de BlackArch}

Afin de garde les choses simples, nos PKGBUILDs sont similaires avec ceux du
Arch User Repository, à quelques différences près. Tous les paquets doivent
être membre du groupe \textit{blackarch}. Il est possible qu'un paquet soit
membre de plus d'un groupe.

\subsection{Groupes}

Afin de permettre aux utilisateurs d'installer certaines catégories de paquets
exclusivement, ceux-ci ont étés répertoriés en plusieurs groupes. Ces groupes
permettent aux utilisateurs d'exécuter la commande 
\textit{pacman -S \textless nom\_du\_groupe \textgreater } afin d'installer tous
 les paquets compris dans ce groupe.

\subsubsection{blackarch}

En théorie, le groupe \textit{blackarch} doit contenir tous les paquets qu'offre
BlackArch. La majorité des paquets devraient être membre de ce groupe.

\subsubsection{blackarch-anti-forensic}

Groupe contenant les paquets utilisés pour contrer les tentatives 
d'investigation informatique de systèmes. Vous y trouverez des paquets offrants
des solutions de cryptage de données, stéganographie, modification d'attributs 
de fichiers, et tout autre type de logiciels permettant la modification d'un
système de fichiers dans le but d'y masquer de l'information.

Exemples: luks, TrueCrypt, Timestomp, dd, ropeadope, secure-delete

\subsubsection{blackarch-automation}

Groupe contenant les paquets utilisés pour l'automation de plusieurs tâches.
Puisque les paquets présents dans ce groupe varient énormément en 
fonctionnalités, il vous est suggéré d'y consulter directement les paquets
afin de vous informer sur leurs capacités.

Exemples: blueranger, tiger, wiffy

\subsubsection{blackarch-backdoor}


Groupe contenant les paquets reliés à l'exploitation ou à l'ouverture de 
backdoors sur de systèmes vulnérables.

Exemples: backdoor-factory, rrs, weevely

\subsubsection{blackarch-binary}

Groupe contenant les paquets reliés aux fichiers binaires. Vous y trouverez des
paquets permettant l'extraction d'archives d'un firmware, des logiciels de
rétro-ingénierie, logiciels aidant à l'exploitation de fichiers binaires, et
autres.

Exemples: binwally, packerid, hex2bin, binwalk

\subsubsection{blackarch-bluetooth}

Groupe contenant les paquets aidant à l'exploitation du standard Bluetooth
(802.15.1).

Exemples: ubertooth, tbear, redfang

\subsubsection{blackarch-code-audit}

Groupe contenant les paquets à fin d'audit de code source. Ces paquets analysent
statiquement le code source de votre projet afin d'y déceler des vulnérabilités.

Exemples: flawfinder, pscan

\subsubsection{blackarch-cracker}

Groupe contenant les paquets utilisés afin de tester la robustesses de multiples
fonctions cryptographiques.

Exemples: hashcat, john, crunch

\subsubsection{blackarch-crypto}

Groupe contenant tous types de paquets ayant rapport de proche ou de loin à la
cryptographie, que ce soit de l'analyse de ou au craquage de hache.

Exemples: ciphertest, xortool, sbd

\subsubsection{blackarch-database}

Groupe contenant les paquets reliés à l'exploitation de bases de données.

Exemples: metacoretex, blindsql

\subsubsection{blackarch-debugger}

Groupe contenant les paquets permettant à l'usager d’investiguer le
fonctionnement interne d'un programme lors de son utilisation.

Exemples: radare2, shellnoob

\subsubsection{blackarch-decompiler}

Groupe contenant les paquets qui tentent d'effectuer de la rétro-ingénierie sur
un programme compilé.

Exemples: flasm, jd-gui

\subsubsection{blackarch-defensive}

Groupe contenant les paquets utilisés afin de protéger un système des logiciels
et usagers malveillants.

Exemples: arpon, chkrootkit, sniffjoke

\subsubsection{blackarch-disassembler}

Ce groupe est fortement similaire à \textit{blackarch-decompiler} dans le sens
que ces deux groupes tentent d'effectuer une rétro-ingénierie générique de
fichiers exécutables, mais les paquets compris dans
\textit{blackarch-disassembler} en extrairont le code assembleur à la place d'un
pseudocode.

Exemples: inguma, radare2

\subsubsection{blackarch-dos}

Groupe contenant les paquets aidant aux attaques de types \textit{déni de service}.

Exemples: 42zip, nkiller2

\subsubsection{blackarch-drone}

Groupe contenant les paquets aidant à l'entretien et l'utilisation de drones.

Exemples: meshdeck, skyjack

\subsubsection{blackarch-exploitation}

Groupe contenant les paquets utilisés lors de l'exploitation de multiples 
programmes et services.

Exemples: armitage, metasploit, zarp

\subsubsection{blackarch-fingerprint}

Groupe contenant les paquets utilisés à fin de reconnaissance réseautique.

Exemples: dns-map, p0f, httprint

\subsubsection{blackarch-firmware}

Groupe contenant les paquets utilisés lors de l'exploitation de firmware.

Exemples: None yet, amend asap.

\subsubsection{blackarch-forensic}

Groupe contenant les paquets utilisés lors d'investigation informatique portant
sur le recouvrement de données potentiellement effacés.

Exemples: aesfix, nfex, wyd

\subsubsection{blackarch-fuzzer}

Groupe contenant les paquets aidant au fuzzing. 

Exemples: msf, mdk3, wfuzz

\subsubsection{blackarch-hardware}

Groupe contenant les paquets aidant à l'exploitation de matériel physique.

Exemples: arduino, smali

\subsubsection{blackarch-honeypot}

Groupe contenant les paquets agissant comme des \textit{honeypots}, c'est-à-dire
des services sécuritaire, mais qui laissent croire à un usager malveillant
autrement afin d'observer son comportement.

Exemples: artillery, bluepot, wifi-honey

\subsubsection{blackarch-keylogger}

Groupe contenant les paquets utilisés lors d'installation de \textit{keylogging}
, soit l'enregistrement de touches clavier sur un système qui pourrait être
utilisé par un usager tiers afin d'en soutirer des informations confidentielles.

Exemples: klogger, logkeys, xspy

\subsubsection{blackarch-malware}

Groupe contenant tous types de paquets ayant un lien au \textit{malware}, que ce
soit à l'utilisation de ceux-ci ou à leur détection.

Exemples: malwaredetect, peepdf, yara

\subsubsection{blackarch-misc}

Groupe contenant les paquets ne semblant pas prendre part dans aucunes autres
catégories.

Exemples: oh-my-zsh-git, winexe, stompy

\subsubsection{blackarch-mobile}

Groupe contenant les paquets aidant aux activités d'entretien et d'audits 
d'applications, de développement d'application, et de matériel mobile..

Exemples: android-sdk-platform-tools, android-udev-rules

\subsubsection{blackarch-networking}

Groupe contenant les paquets aidant à multiples tâches reliés à la réseautique.

Exemples: Anything pretty much

\subsubsection{blackarch-nfc}

Groupe contenant les paquets utilisant la technologie \textit{NFC} (near-field
communications).

Exemples: nfcutils

\subsubsection{blackarch-packer}

Groupe contenant les paquets utilisant de loin ou de près les \textit{packers},
soit des logiciels qui créent des applications contenant un programme 
malveillant caché à l'intérieur.

Exemples: packerid

\subsubsection{blackarch-proxy}

Groupe contenant les paquets agissant tel un proxy, soit un logiciel redirigeant
le trafic internet afin de l'analyser ou de le modifier.

Exemples: burpsuite, ratproxy, sslnuke

\subsubsection{blackarch-recon}

Groupe contenant les paquets servant à la recherche de services exploitable dans
le monde réel.

Exemples: canri, dnsrecon, netmask

\subsubsection{blackarch-reversing}

Groupe contenant les paquets de décompilation, de désassemblage, et tout autres
types de programmes ayant à voir avec la rétro-ingénierie. 

Exemples: capstone, radare2, zerowine

\subsubsection{blackarch-scanner}

Groupe contenant les paquets utilisés afin de scanner multiples systèmes afin
d'y déceler des vulnérabilités actives.

Exemples: scanssh, tiger, zmap

\subsubsection{blackarch-sniffer}

Groupes contenant les paquets aidant dans l'analyse de trafic réseau.

Exemples: hexinject, pytactle, xspy

\subsubsection{blackarch-social}

Groupe contenant les paquets attaquant principalement les sites de réseaux
sociaux.

Exemples: jigsaw, websploit

\subsubsection{blackarch-spoof}

Groupe contenant les paquets qui tentent de dissimuler un attaquant, de tel que
la victime ne puisse le percevoir comme un acteur malicieux.

Exemples: arpoison, lans, netcommander

\subsubsection{blackarch-threat-model}

Groupe contenant les paquets utilisé pour enregistré et créer des rapports sur
plusieurs modèles de menaces pouvant se produire lors d'un scénario spécifique.

Exemples: magictree

\subsubsection{blackarch-tunnel}

Groupe contenant les paquets utilisés afin d'encapsuler le contenu réseau via un
tiers parti, afin de dissimuler la provenance initiale de données. 

Exemples: ctunnel, iodine, ptunnel

\subsubsection{blackarch-unpacker}

Groupe contenant les paquets utilisés lors de l'extraction de logiciel
malveillants de possibles logiciels \textit{packés} ou obfuscé.

Exemples: js-beautify

\subsubsection{blackarch-voip}

Groupe contenant les paquets utilisés afin d'effectuer des actions ou l'analyse
de la téléphonie VoIP.

Exemples: iaxflood, rtp-flood, teardown

\subsubsection{blackarch-webapp}

Groupe contenant les paquets effectuant des opérations de validation de sécurité
sur des applications accessible via l'internet, afin d'en garantir leurs 
robustesse. 

Exemples: metoscan, whatweb, zaproxy

\subsubsection{blackarch-windows}

Groupe contenant les paquets Windows qui s'exécuteront via Wine.

Exemples: 3proxy-win32, pwdump, winexe

\subsubsection{blackarch-wireless}

Groupe contenant les paquets utilisés lors d'audits de sécurité sur une 
infrastructure sans-fil.

Exemples: airpwn, mdk3, wiffy

\section{Structure du dépôt}

Vous pouvez trouver le dépot principal de BlackArch à 
\href{https://github.com/BlackArch/blackarch}
{https://github.com/BlackArch/blackarch}. Il existe aussi plusieurs dépots de
code à \href{https://github.com/BlackArch}{https://github.com/BlackArch}.

À l'intérieur du dépôt principal, vous trouverez les trois dossiers principaux,
soit :

\begin{itemize}
\item docs - Contient la documentation du projet
\item packages - Contient les fichiers PKGBUILD
\item scripts - Contient multiples scripts d'automation
\end{itemize}

\subsection{Scripts}

Voici une référence des scripts que vous trouverez dans le dossier
\verb|scripts/| :

\begin{itemize}
\item baaur - Éventuellement, ce script enverra nos paquets à la banque de donné
du AUR.
\item babuild - Construit un paquet.
\item bachroot - Prépare un chroot afin d'effectuer des tests.
\item baclean - Néttoie les enciens fichiers .pkg.tar.xz du dépôt.
\item baconflict - Éventuellement, ce script remplacera \verb|scripts/conflicts|.
\item bad-files - Trouve des fichiers erronés dans les paquets construits.
\item balock - Crée ou termine le lock sur le dépôt de paquets.
\item banotify - Indique IRC des push de paquets.
\item barelease - Envoie les paquets au dépôt.
\item baright - Affiche les informations de copyright de BlackArch.
\item basign - Signe les paquets.
\item basign-key - Signe une clé.
\item blackman - Opère sensiblement comme pacman mais construits les paquets à
partir de git. (à ne pas confondre avec le programme Blackman, écrit par nzr)
\item check-groups - Valide les groupes.
\item checkpkgs - Valide les paquets.
\item conflicts - Valide les paquets pour des conflits.
\item dbmod - Modifie la base de donnée d'un paquet.
\item depth-list - Crée une liste de dépendences ordonnée par profondeur.
\item deptree - Crée une arbre de dépendance, ne listant que les paquets obtenu
via BlackArch.
\item get-blackarch-deps - Affiche une liste des dépendances de paquets obtenu 
via BlackArch pour un paquet donné.
\item get-official - Obtien les paquets officiels.
\item list-loose-packages - Liste les paquets qui ne se trouvent dans aucun
groupes et ne sont pas des dépendances pour d'autres paquets.
\item list-needed - Liste les dépendances manquantes.
\item list-removed - Liste les paquets qui sont dans le dépôt mais pas dans le
répertoire git.
\item list-tools - Liste les outils.
\item outdated - Cherche les paquets qui ne sont plus à jour lorsque comparés 
avec le répertoire git.
\item pkgmod - Modifie un paquets de construction.
\item pkgrel - Incrémente le pkgrel d'un paquet.
\item prep - Nettoie un fichier PKGBUILD de ses erreurs de style et autre.
\item sitesync - Synchronise entre une copie locale et sa version distante.
\item size-hunt - Cherche de gros paquets.
\item source-backup - Sauvegarde les fichiers sources d'un paquets.
\end{itemize}

\section{Contribuer au dépôt}
Cette section indique comment contribuer au projet BlackArch Linux. Nous
acceptons les pull requests de toutes tailles, que ce soit de petites
modifications ou un ajout de paquet.\\Pour recevoir de l'aide, n'hésitez pas à
nous contacter.
\\\\
Tout le monde est bienvenue à contrivuer. Toute contribution sera appréciée.

\subsection{Tutoriels requis.}
Veuillez lire les documents suivants avant de contribuer:
\begin{itemize}
\item \href{http://wiki.archlinux.fr/Standard_paquetage}{Arch Wiki FR:
Standard paquetage}
\item \href{http://wiki.archlinux.fr/Makepkg}{Arch Wiki FR: Makepkg}
\item \href{http://wiki.archlinux.fr/PKGBUILD}{Arch Wiki FR: PKGBUILD}
\end{itemize}

\subsection{Étapes à suivre afin de contribuer}
Afin de contribuer au projet BlackArch Linux, suivez ces étapes:
\begin{enumerate}
\item Forkez le dépôt sur github à partir de
\url{https://github.com/BlackArchLinux/blackarchlinux}.
\item Effectuez vos modifications sur les fichiers nécéssaires (e.g. PKGBUILD,
.patch files, etc).
\item Commitez vos changements.
\item Poussez vos changements.
\item Créez un pull request vers le dépôt.
\end{enumerate}

\subsection{Exemple}
L'exemple suivant démontre la soumissions d'un nouveau paquets au projet
BlackArch. Nous utilisons \href{https://wiki.archlinux.org/index.php/yaourt}{yaourt}
(vous pouvez aussi utiliser pacaur) afin d'obtenir le PKGBUILD existant de
\textbf{nfsshell}, et nous pouvons ensuite l'éditer, au besoin.

\subsubsection{Obtenir le PKGBUILD}
En utilisant yaourt ou pacaur, obtenir \textit{PKGBUILD}:
{\small
\color{gray}
\begin{verbatim}
user@blackarchlinux $ yaourt -G nfsshell
==> Download nfsshell sources
x LICENSE
x PKGBUILD
x gcc.patch
user@blackarchlinux $ cd nfsshell/
\end{verbatim}
}

\subsubsection{Nettoiement du PKGBUILD}
Nettoyer le fichier \textit{PKGBUILD} à l'aide de \textit{prep}, afin de sauver
du temps:
{\small
\color{gray}
\begin{verbatim}
user@blackarchlinux nfsshell $ ./blarckarch/scripts/prep PKGBUILD
cleaning 'PKGBUILD'...
expanding tabs...
removing vim modeline...
removing id comment...
removing contributor and maintainer comments...
squeezing extra blank lines...
removing '|| return'...
removing leading blank line...
removing $pkgname...
removing trailing whitespace...
\end{verbatim}
}

\subsubsection{Ajustement du PKGBUILD}
Ajustez le fichier \textit{PKGBUILD}:
{\small
\color{gray}
\begin{verbatim}
user@blackarchlinux nfsshell $ vi PKGBUILD
\end{verbatim}
}

\subsubsection{Construire le paquet}
Construire le paquet:
{\small
\color{gray}
\begin{verbatim}
user@blackarchlinux nfsshell $ makepkg -sf
==> Making package: nfsshell 19980519-1 (Mon Dec 2 17:23:51 CET 2013)
==> Checking runtime dependencies...
==> Checking buildtime dependencies...
==> Retrieving sources...
-> Downloading nfsshell.tar.gz...
% Total % Received % Xferd Average Speed Time Time Time
CurrentDload Upload Total Spent Left Speed100 29213 100 29213 0
0 48150 0 --:--:-- --:--:-- --:--:-- 48206
-> Found gcc.patch
-> Found LICENSE
...
<beaucoup d'output du processus de build et compilation ici>
...
==> Leaving fakeroot environment.
==> Finished making: nfsshell 19980519-1 (Mon Dec 2 17:23:53 CET 2013)
\end{verbatim}
}

\subsubsection{Installation et test du paquet}
Installation et test du paquet:
{\small
\color{gray}
\begin{verbatim}
user@blackarchlinux nfsshell $ pacman -U nfsshell-19980519-1-x86_64.pkg.tar.xz
user@blackarchlinux nfsshell $ nfsshell # test it
\end{verbatim}
}

\subsubsection{Ajout, commit et push du paquet}
Ajoute, commit et push du paquet:
{\small
\color{gray}
\begin{verbatim}
user@blackarchlinux nfsshell $ cd /blackarchlinux/packages
user@blackarchlinux ~/blackarchlinux/packages $ mv ~/nfsshell .
user@blackarchlinux ~/blackarchlinux/packages $ git commit -am nfsshell && git push
\end{verbatim}
}

\subsubsection{Création d'un pull request}
Créez un pull request de votre fork vers
 \href{https://github.com/BlackArchLinux/blackarchlinux}{le projet officiel}:
{\small
\color{gray}
\begin{verbatim}
firefox https://github.com/<contributor>/blackarchlinux
\end{verbatim}
}

\subsubsection{Ajout d'un remote pour l'upstream}
Une fois votre modification faite, il serait une bonne idée d'ajouter le projet
officiel comme remote, afin d'ajouter les modifications au projet à votre fork.
Cela vous permettra de synchroniser votre fork avec BlackArch.
{\small
\color{gray}
\begin{verbatim}
user@blackarchlinux ~/blackarchlinux $ git remote -v
origin <the url of your fork> (fetch)
origin <the url of your fork> (push)
user@blackarchlinux ~/blackarchlinux $ git remote add upstream https://github.com/blackarch/blackarch
user@blackarchlinux ~/blackarchlinux $ git remote -v
origin <the url of your fork> (fetch)
origin <the url of your fork> (push)
upstream https://github.com/blackarch/blackarch (fetch)
upstream https://github.com/blackarch/blackarch (push)
\end{verbatim}
}
Par défaut, git devrait push directement à l'origine, mais assurez-vous que
votre configuration git est correcte. Cela ne sera pas un problème à moins
d'avoir les droits de commit puisque vou sne pourrez pas push upstream sans
cela.

Si vous avez les droits de commit, vous aurez possiblement plus de succès en
utilisant \textit{git@github.com:blackarch/blackarch.git}. Le choix est votre.

\subsection{Requêtes}
\begin{enumerate}
\item N'ajoutez pas de commentaire \textbf{Maintainer} ou \textbf{Contributor}
au fichiers \textit{PKGBUILD}. Ajoutez ces commentaires à la section
\textbf{AUTHORS} du guide BlackArch.
\item Afin de conserver la consistance du projet, veuillez suivre les règles
générales de style des autres fichiers \textit{PKGBUILD} en utilisant un
indentation de deux espaces.
\end{enumerate}

\subsection{Conseils généraux}
\href{http://wiki.archlinux.org/index.php/Namcap}{namcap} peut trouver des
erreurs dans vos paquets. Il est recommendé de l'utiliser avant de soumettre vos
modifications.

%%% APPENDIX %%%
\appendix
%%%%%%%%%%%%%%%%%%%%%%%%%%%%%%%%%%%%%%%%%%%%%%%%%%%%%%%%%%%%%%%%%%%%%%%%%%%%%%%%
%                                                                              %
% BlackArch Linux appendix template                                            %
%                                                                              %
%%%%%%%%%%%%%%%%%%%%%%%%%%%%%%%%%%%%%%%%%%%%%%%%%%%%%%%%%%%%%%%%%%%%%%%%%%%%%%%%

\appendix

\chapter{Appendix}

\section{FAQs}

\section{AUTHORS}
The following people have contributed directly to BlackArch:
\begin{itemize}
\item Evan Teitelman <teitelmanevan at gmail dot com>
\item Tyler Bennnett <tylerb at trix2voip dot com>
\item Levon Kayan <noptrix at nullsecurity dot net>
\item Jeremy Lynch <jl at adminempire dot com>
\item Ari Mizrahi <codemunchies at gmail dot com>
\item fnord0 <fnord0 at riseup dot net>
\item nrz <nrz at nullsecurity dot net>
\item Ellis Kenyo <elken.tdos at gmail dot com>
\item CaledoniaProject <the.warl0ck.1989 at gmail dot com>
\item sudokode <sudokode at gmail dot com>
\item Valentin Churavy <v.churavy at gmail dot com>
\item Boy Sandy Gladies Arriezona <reno.esper at gmail dot com>
\item Mathias Nyman <None>
\item Johannes Löthberg <demizide at gmail dot com>
\end{itemize}
The following people contributed directly to ArchPwn, which has been merged
into BlackArch:
\begin{itemize}
\item Francesco Piccinno <stack.box at gmail dot com>
\item jensp <jens at jenux.homelinux dot org>
\item Valentin Churavy <v.churavy at gmail dot com>
\end{itemize}
We have taken build code from the following people:
\begin{itemize}
\item  3ED <krzysztof1987 at gmail dot com>
\item  AUR Perl <aurperl at juster dot info>
\item  Aaron Griffin <aaron at archlinux dot org>
\item  Abakus <java5 at arcor dot de>
\item  Adam Wolk <netprobe at gmail dot com>
\item  Aleix Pol <aleixpol at kde dot org>
\item  Aleshus <aleshusi at gmail dot com>
\item  Alessandro Pazzaglia <jackdroido at gmail dot com>
\item  Alessandro Sagratini <ale\_sagra at hotmail dot com>
\item  Alex Cartwright <alexc223 at googlemail dot com>
\item  Alexander De Sousa <archaur.xandy21 at spamgourmet dot com>
\item  Alexander Rødseth <rodseth at gmail dot com>
\item  Allan McRae <allan at archlinux dot org>
\item  AmaN <gabroo.punjab.da at gmail dot com>
\item  Andre Klitzing <aklitzing at online dot de>
\item  Andrea Scarpino <andrea at archlinux dot org>
\item  Andreas Schönfelder <passtschu at freenet dot de>
\item  Andrej Gelenberg <andrej.gelenberg at udo dot edu>
\item  Angel Velasquez <angvp at archlinux dot org>
\item  Antoine Lubineau <antoine at lubignon dot info>
\item  Anton Bazhenov <anton.bazhenov at gmail>
\item  Arkham <arkham at archlinux dot us>
\item  Arthur Danskin <arthurdanskin at gmail dot com>
\item  Balda Baldanos \_at\_ gmail dot com
\item  Balló György <ballogyor+arch at gmail dot com>
\item  Bartek Piotrowski <barthalion at gmail dot com>
\item  Bartosz Feński <fenio at debian dot org>
\item  Bartłomiej Piotrowski <nospam at bpiotrowski dot pl>
\item  Bogdan Szczurek <thebodzio at gmail dot com>
\item  Brad Fanella <bradfanella at archlinux dot us>
\item  Brian Bidulock <bidulock at openss7 dot org>
\item  C Anthony Risinger <anthony at xtfx dot me>
\item  CRT <crt.011 at gmail dot com>
\item  Can Celasun <dcelasun at gmail dot com>
\item  Chaniyth <chaniyth at yahoo dot com>
\item  Chris Brannon <cmbrannon79 at gmail dot com>
\item  Chris Giles <Chris.G.27 at gmail dot com> \& daschu117
\item  Christoph Siegenthaler <csi at gmx dot ch>
\item  Christoph Zeiler <archNOSPAM at moonblade dot org>
\item  Clément DEMOULINS <clement at archivel dot fr>
\item  Corrado Primier <bardo at aur.archlinux dot org>
\item  Daenyth <Daenyth+Arch at gmail dot com>
\item  Dale Blount <dale at archlinux dot org>
\item  Damir Perisa <damir.perisa at bluewin dot ch>
\item  Dan Fuhry <dan at fuhry dot us>
\item  Dan Serban <dserban01 at yahoo dot com>
\item  Daniel A. Campoverde Carrión
\item  Daniel Golle
\item  Daniel Griffiths <ghost1227 at archlinux dot us>
\item  Daniel J Griffiths <ghost1227 at archlinux dot us>
\item  Daniel Micay <danielmicay at gmail dot com>
\item  Dave Reisner <dreisner at archlinux dot org>
\item  Dawid Wrobel <cromo at klej dot net>
\item  Devaev Maxim <mdevaev at gmail dot com>
\item  Devin Cofer <ranguvar at archlinux dot us>
\item  DigitalPathogen <aur at InfoSecResearchLabs dot co dot uk>
\item  DigitalPathogen <aur at digitalpathogen dot co dot uk>
\item  Dmitry A. Ilyashevich <dmitry.ilyashevich at gmail dot com>
\item  Dominik Heidler <dheidler at gmail dot com>
\item  DrZaius <lou at fakeoutdoorsman dot com>
\item  Ebubekir KARUL <ebubekirkarul at yandex dot com>
\item  Eduard "bekks" Warkentin <eduard.warkentin at gmail dot com>
\item  Elmo Todurov <todurov at gmail dot com>
\item  Emmanuel Gil Peyrot <linkmauve at linkmauve dot fr>
\item  Eric Belanger <eric at archlinux dot org>
\item  Ermak <ermak at email dot it>
\item  Evangelos Foutras <evangelos at foutrelis dot com>
\item  Fabian Melters <melters at gmail dot com>
\item  Fabiano Furtado <fusca14 at gmail dot com>
\item  Federico Quagliata (quaqo) <linux at quaqo dot org>
\item  Firmicus <francois.archlinux dot org>
\item  Florian Pritz <bluewind at jabber dot ccc dot de>
\item  Florian Pritz <flo at xinu dot at>
\item  Francesco Piccinno <stack.box at gmail dot com>
\item  François Charette <francois at archlinux dor org>
\item  Gaetan Bisson <bisson at archlinux dot org>
\item  Geoffroy Carrier <geoffroy.carrier at koon.fr>
\item  Georg Grabler (STiAT) <ggrabler at gmail dot com>
\item  George Hilliard <gh403 at msstate dot edu>
\item  Gerardo Exequiel Pozzi <vmlinuz386 at yahoo dot com dot ar>
\item  Gilles CHAUVIN <gcnweb at gmail dot com>
\item  Giovanni Scafora <giovanni at archlinux dot org>
\item  Gordin <9ordin  at t gmail dot com>
\item  Guillaume ALAUX <guillaume at archlinux dot org>
\item  Guillermo Ramos <0xwille at gmail dot com>
\item  Gustavo Alvarez <sl1pkn07 at gmail dot com>
\item  Hugo Doria <hugo at archlinux dot org>
\item  Hyacinthe Cartiaux <hyacinthe.cartiaux at free dot fr>
\item  James Fryman <jfryman at gmail dot com>
\item  Jan "heftig" Steffens <jan.steffens at gmail dot com>
\item  Jan de Groot <jgc at archlinux dot org>
\item  Jaroslav Lichtblau <dragonlord at aur dot archlinux dot org>
\item  Jaroslaw Swierczynski <swiergot at aur dot archlinux dot org>
\item  Jason Chu <jason at archlinux dot org>
\item  Jason R Begley (jayray at digitalgoat dot com>
\item  Jason Rodriguez <jason\item aur at catloaf dot net>
\item  Jason St. John <jstjohn at purdue dot edu>
\item  Jawmare <victor2008 at gmail dot com>
\item  Jeff Mickey <jeff at archlinux dot org>
\item  Jens Pranaitis <jens at chaox dot net>
\item  Jens Pranaitis <jens at jenux dot homelinux dot org>
\item  Jinx <jinxware at gmail dot com>
\item  John D Jones III <j[nospace]n[nospace]b[nospace]e[nospace]k[nospace]1972 at gmail dot com>
\item  John Proctor <jproctor at prium dot net>
\item  Jon Bergli Heier <snakebite at jvnv dot net>
\item  Jonas Heinrich <onny at project\item insanity dot org>
\item  Jonathan Steel <jsteel at aur dot archlinux dot org>
\item  Joris Steyn <jorissteyn at gmail dot com>
\item  Josh VanderLinden <arch at cloudlery dot com>
\item  Jozef Riha <jose1711 at gmail dot com>
\item  Judd Vinet <jvinet at zeroflux dot org>
\item  Juergen Hoetzel <jason at archlinux dot org>
\item  Juergen Hoetzel <juergen at archlinux dot org>
\item  Justin Davis <jrcd83 at gmail dot com>
\item  Kaiting Chen <kaitocracy at gmail dot com>
\item  Kaos < gianlucaatlas at gmail dot com >
\item  Kevin Piche <kevin at archlinux dot org>
\item  Kory Woods <kory \_at\_ virlo >dot< net>
\item  Kyle Keen <keenerd at gmail dot com>
\item  Larry Hajali <larryhaja at gmail dot com>
\item  LeCrayonVert < greenarrow at archlinux dot us>
\item  Le\_suisse <lesuisse dot dev+aur at gmail dot com>
\item  Lekensteyn <lekensteyn at gmail dot com>
\item  Limao Luo <luolimao+AUR at gmail dot com>
\item  Lucien Immink <l.immink at student dot fnt dot hvu dot nl>
\item  Lukas Fleischer <archlinux at cryptocrack dot de>
\item  Manolis Tzanidakis
\item  Marcin "avalan" Falkiewicz <avalatron at gmail dot com>
\item  Mariano Verdu <verdumariano at gmail dot com>
\item  Marti Raudsepp <marti at juffo dot org>
\item  MatToufoutu <mattoufootu at gmail dot com>
\item  Matthew Sharpe <matt.sharpe at gmail dot com>
\item  Mauro Andreolini <mauro.andreolini at unimore dot it>
\item  Max Pray a.k.a. Synthead <synthead at gmail dot com>
\item  Max Roder <maxroder at web dot de>
\item  Maxwell Pray a.k.a. Synthead <synthead at gmail dot com>
\item  Maxwell Pray a.k.a. Synthead <synthead1 at gmail dot com>
\item  Mech <tiago.bmp at gmail dot com>
\item  Michael Düll <mail at akurei dot me>
\item  Michael P <ptchinster at archlinux dot us>
\item  Michal Krenek <mikos at sg1 dot cz>
\item  Michal Zalewski <lcamtuf at coredumpdotcx>
\item  Miguel Paolino <mpaolino at gmail dot com>
\item  Miguel Revilla <yo at miguelrevilla dot com>
\item  Mike Roberts <noodlesgc at gmail dot com>
\item  Mike Sampson <mike at sambodata dot com>
\item  Nassim Kacha <nassim.kacha at gmail dot com>
\item  Nicolas Pouillard <nicolas.pouillard at gmail dot com>
\item  Nicolas Pouillard https://nicolaspouillard.fr
\item  Niklas Schmuecker (IRC: nisc) <nschmuecker gmail dot com>
\item  Oleander Reis <oleander at oleander dot cc>
\item  Olivier Le Moal <mail at olivierlemoal dot fr>
\item  Olivier Médoc "oliv" <o\_medoc at yahoo dot fr>
\item  Pascal E. <archlinux at hardfalcon dot net>
\item  Patrick Leslie Polzer <leslie.polzer at gmx dot net>
\item  Paul Mattal <paul at archlinux dot org>
\item  Paul Mattal <pjmattal at elys dot com>
\item  Pengyu CHEN <cpy.prefers.you at gmail dot com>
\item  Peter Wu <lekensteyn at gmail dot com>
\item  Philipp 'TamCore' B. <philipp at tamcore dot eu>
\item  Pierre Schmitz <pierre at archlinux dot de>
\item  Pranay Kanwar <pranay.kanwar at gmail dot com>
\item  Pranay Kanwar <warl0ck at metaeye dot org>
\item  PyroPeter <googlemail dot com at abi1789>
\item  PyroPeter <googlemail.com at abi1789>
\item  Ray Rashif <schiv at archlinux dot org>
\item  Remi Gacogne <rgacogne\item arch at coredump dot fr>
\item  Renan Fernandes <renan at kauamanga dot com dot br>
\item  Richard Murri <admin at richardmurri dot com>
\item  Roberto Alsina <ralsina at kde dot org>
\item  Robson Peixoto <robsonpeixoto at gmail dot com>
\item  Roel Blaauwgeers <roel at ttys0 dot nl>
\item  Rorschach <r0rschach at lavabit dot com>
\item  Ruben Schuller <shiml at orgizm dot net>
\item  Rudy Matela <rudy at matela dot com dot br>
\item  Ryon Sherman <ryon.sherman at gmail dot com>
\item  Sabart Otto \item  Seberm <seberm at gmail dot com>
\item  SakalisC <chrissakalis at gmail dot com>
\item  Sam Stuewe <halosghost at archlinux dot info>
\item  SanskritFritz <SanskritFritz at gmail dot com)
\item  Sarah Hay <sarahhay at mb dot sympatico dot ca>
\item  Sebastian Benvenuti <sebastianbenvenuti at gmail dot com>
\item  Sebastian Nowicki <sebnow at gmail dot com>
\item  Sebastien Duquette <ekse.0x at gmail dot com>
\item  Sebastien LEDUC <sebastien at sleduc dot fr>
\item  Sebastien Leduc <sebastien at sleduc dot fr>
\item  Sergej Pupykin <pupykin.s+arch at gmail dot com>
\item  Sergio Rubio <rubiojr at biondofu dot net>
\item  Sheng Yu <magicfish1990 at gmail dot com>
\item  Simon Busch <morphis at gravedo dot de>
\item  Simon Legner <Simon.Legner at gmail dot com>
\item  Sirat18 <aur at sirat18 dot de>
\item  SpepS <dreamspepser at yahoo dot it>
\item  Spider.007 <archPackage at spider007 dot net>
\item  Stefan Seering
\item  Stephane Travostino <stephane.travostino at gmail dot com>
\item  Stéphane Gaudreault <stephane at archlinux dot org>
\item  Sven Kauber <celeon at gmail dot com>
\item  Sven Schulz <omee at archlinux dot de>
\item  Sébastien Duquette <ekse.0x at gmail dot com>
\item  Sébastien Luttringer <seblu at archlinux dot org>
\item  TDY <tdy at gmx dot com>
\item  Teemu Rytilahti <tpr at iki dot fi>
\item  Testuser\_01 <mail at nico\item siebler dot de>
\item  Thanx <thanxm at gmail dot com>
\item  Thayer Williams <thayer at archlinux dot org>
\item  Thomas S Hatch <thatch45 at gmail dot com>
\item  Thorsten Töpper <atsutane\item aur at freethoughts dot de>
\item  Thorsten Töpper <atsutane\item tu at freethoughts dot de>
\item  Tilmann Becker <tilmann.becker at freenet dot de>
\item  Timothy Redaelli <timothy.redaelli at gmail dot com>
\item  Timothée Ravier <tim at siosm dot fr>
\item  Tino Reichardt
\item  Tobias Kieslich <tobias at justdreams dot de>
\item  Tobias Powalowski <tpowa at archlinux dot org>
\item  Tom K <tomk at runbox dot com>
\item  Tom Newsom <Jeepster at gmx dot co dot uk>
\item  Tomas Lindquist Olsen <tomas at famolsen dot dk>
\item  Travis Willard <travisw at wmpub dot ca>
\item  Valentin Churavy <v.churavy at gmail dot com>
\item  ViNS <gladiator at fastwebnet dot it>
\item  Vlatko Kosturjak <kost at linux dot hr>
\item  Wes Brown <wesbrown18 at gmail dot com>
\item  William Rea <sillywilly at gmail dot com>
\item  Xavier Devlamynck <magicrhesus at ouranos dot be>
\item  Xiao\item Long Chen <chenxiaolong at cxl dot epac dot to>
\item  aeolist <aeolist at hotmail dot com>
\item  ality at pbrane dot org
\item  astaroth <astaroth\_ at web dot de>
\item  bender02 at archlinux dot us
\item  billycongo <billycongo at gmail dot com>
\item  bslackr <brendan at vastactive dot com>
\item  cbreaker <cbreaker at tlen dot pl>
\item  chimeracoder <dev@chimeracoder.net>
\item  damir <damir at archlinux.org>
\item  danitool
\item  darkapex <me at jailuthra dot in>
\item  daronin
\item  dkaylor <dpkaylor at gmail dot com>
\item  dobo <dobo90\_at\_gmail dot com>
\item  dorphell <dorphell at archlinux dot org>
\item  evr <evanroman  at  gmail>
\item  fnord0 <fnord0 at riseup dot net>
\item  fxbru <frxbru at gmail>
\item  hcar
\item  icarus <icarus.roaming at gmail dot com>
\item  iceman <icemanf at gmail dot com>
\item  kastor <kastor at fobos dot org dot ar>
\item  kfgz <kfgz at interia dot pl>
\item  linuxSEAT <linuxSEAT at gmail dot com>
\item  m4xm4n <max at maxfierke dot com>
\item  mar77i <mysatyre at gmail dot com>
\item  marc0s <marc0s at fsfe dot org>
\item  mickael9 <mickael9 at gmail dot com>
\item  nblock <nblock at archlinux dot us>
\item  nofxx <x at nofxx dot com>
\item  onny <onny at project
\item  pootzko <pootzko at gmail dot com>
\item  revel <revel at muub dot net>
\item  rich\_o <rich\_o at lavabit dot com>
\item  s1gma <s1gma at mindslicer dot com>
\item  sandman <r.coded at gmail dot com>
\item  sebikul <sebikul at gmail dot com>
\item  sh0 <mee at sh0 dot org>
\item  shild <sxp at bk dot ru>
\item  simo <simo at archlinux dot org>
\item  snuo
\item  sudokode <sudokode at gmail dot com>
\item  tobias <tobias at archlinux dot org>
\item  trashstar <trash at ps3zone dot org>
\item  unexist <unexist at subforge dot org>
\item  untitled <rnd0x00 at gmail dot com>
\item  virtuemood <virtue at archlinux dot us>
\item  wido <widomaker2k7 at gmail dot com>
\item  wodim <neikokz at gmail dot com>
\item  yannsen <ynnsen at gmail dot com>
\end{itemize}


\end{document}

%%% EOF %%%